% singular_cases.tex: how a target voxel m (blue) can sit relative to a source
% voxel n (red), with the shared corner, edge or face where R = 0 highlighted.
% Uses the palette and styles defined in preamble.tex.
\begin{tikzpicture}[font=\small]
  \def\W{2.95}   % panel width
  \def\G{0.375} % gap
  \def\H{2.8}   % panel height
  \def\a{0.9}   % voxel edge
  % \cuboid{x0}{y0}{z0}{colour}: translucent axis-aligned cube in oblique projection
  \newcommand\cuboid[4]{%
    \begin{scope}[x={(1cm,0cm)}, y={(0.5cm,0.36cm)}, z={(0cm,1cm)}]
      \draw[black!35, dashed] (#1,#2+\a,#3) -- (#1+\a,#2+\a,#3)
        (#1,#2+\a,#3) -- (#1,#2,#3) (#1,#2+\a,#3) -- (#1,#2+\a,#3+\a);
      \filldraw[fill=#4, fill opacity=0.12, draw=inkcol]
        (#1,#2,#3) -- (#1+\a,#2,#3) -- (#1+\a,#2,#3+\a) -- (#1,#2,#3+\a) -- cycle;
      \filldraw[fill=#4, fill opacity=0.12, draw=inkcol]
        (#1+\a,#2,#3) -- (#1+\a,#2+\a,#3) -- (#1+\a,#2+\a,#3+\a) -- (#1+\a,#2,#3+\a) -- cycle;
      \filldraw[fill=#4, fill opacity=0.12, draw=inkcol]
        (#1,#2,#3+\a) -- (#1+\a,#2,#3+\a) -- (#1+\a,#2+\a,#3+\a) -- (#1,#2+\a,#3+\a) -- cycle;
    \end{scope}}
  \tikzset{oblique/.style={x={(1cm,0cm)}, y={(0.5cm,0.36cm)}, z={(0cm,1cm)}}}

  \foreach \n/\title/\how in {0/{Separated}/{},
                              1/{Corner}/{},
                              2/{Edge}/{},
                              3/{Face}/{},
                              4/{Coincident}/{}} {
    \begin{scope}[shift={(\n*\W+\n*\G,0)}]
      \draw[panel] (0,0) rectangle (\W,\H);
      \pgfmathtruncatemacro\letter{97+\n}
      \node[paneltitle] at (0,\H) {\textbf{(\char\letter)} \title};
      \node[below, inner sep=4pt, font=\footnotesize, inkcol] at (\W/2,0) {\how};
    \end{scope}
  }

  % (a) separated
  \begin{scope}[shift={(0.3,0.35)}]
    \cuboid{0}{0}{0}{matcol}
    \cuboid{1.45}{0}{0.9}{vaccol}
    \node[matcol, font=\footnotesize] at (0.55,1.5) {$\Omega_n$};
    \node[vaccol, font=\footnotesize, anchor=north] at (1.9,0.9) {$\Omega_m$};
  \end{scope}
  % (b) shared corner
  \begin{scope}[shift={(\W+\G+0.12,0.2)}]
    \cuboid{0}{0}{0}{matcol}
    \cuboid{\a}{\a}{\a}{vaccol}
    \fill[accent, oblique] (\a,\a,\a) circle (2pt);
  \end{scope}
  % (c) shared edge
  \begin{scope}[shift={(2*\W+2*\G+0.5,0.35)}]
    \cuboid{0}{0}{0}{matcol}
    \cuboid{\a}{0}{\a}{vaccol}
    \draw[accent, line width=2pt, oblique] (\a,0,\a) -- (\a,\a,\a);
  \end{scope}
  % (d) shared face
  \begin{scope}[shift={(3*\W+3*\G+0.4,0.75)}]
    \cuboid{0}{0}{0}{matcol}
    \fill[accent, opacity=0.45, oblique] (\a,0,0) -- (\a,\a,0) -- (\a,\a,\a) -- (\a,0,\a) -- cycle;
    \cuboid{\a}{0}{0}{vaccol}
  \end{scope}
  % (e) coincident
  \begin{scope}[shift={(4*\W+4*\G+0.85,0.65)}]
    \cuboid{0}{0}{0}{matcol}
    \cuboid{0}{0}{0}{vaccol}
    \node[font=\footnotesize, anchor=north] at (0.45,0) {$\Omega_m = \Omega_n$};
  \end{scope}
\end{tikzpicture}
